\section{Introduction}

Forward selection is still routinely used to choose which covariates enter a
regression model, and the confidence intervals reported afterwards are almost
always the ones the software prints for the selected model. Those intervals
ignore the fact that the model was chosen using the same data, and so they do
not have their nominal coverage. We quantify the problem by simulation and
compare against the simplest available remedy, sample splitting.

\section{Simulation}

We generate $n$ observations on $p = 10$ predictors with AR(1) correlation
$\rho$, where $\beta_1 = 1$, $\beta_2 = -1$, and $\beta_3 = \cdots = \beta_{10}
= 0$. For each replication we select variables by forward selection using AIC
and then form intervals for the selected coefficients two ways:

\begin{itemize}
  \item \emph{Naive}: refit the selected model on all $n$ observations and take
        the usual intervals.
  \item \emph{Sample splitting}: select on a random half of the data and fit the
        selected model on the other half.
\end{itemize}

\section{Results}

Figure~\ref{fig:coverage} shows coverage against $n$ for three correlation
levels. The naive intervals cover well below the nominal rate at every sample
size, and increasing $n$ does not help: selection error does not vanish as the
sample grows, because there are always null covariates that happen to look
useful. Sample splitting is close to nominal throughout, at the cost of using
half the data for each task.

\begin{figure}[ht]
  \centering
  \includegraphics[width=\textwidth]{results/figure1.pdf}
  \caption{Coverage of nominal 95\% confidence intervals for coefficients
    selected by forward selection.}
  \label{fig:coverage}
\end{figure}

Table~\ref{tab:coverage} separates the selected coefficients into those that are
truly non-null and those that are truly zero. The naive undercoverage is driven
almost entirely by the null coefficients, which are selected only when their
sample correlation with the outcome is unusually large --- exactly the situation
in which the resulting interval is centered far from zero.

\begin{table}

\caption{\label{tab:coverage}Coverage of nominal 95\% confidence intervals for coefficients selected by forward selection, by whether the selected coefficient is truly zero.}
\centering
\begin{tabular}[t]{rcccc}
\toprule
\multicolumn{1}{c}{ } & \multicolumn{2}{c}{Non-null} & \multicolumn{2}{c}{Null} \\
\cmidrule(l{3pt}r{3pt}){2-3} \cmidrule(l{3pt}r{3pt}){4-5}
$n$ & Naive & Sample splitting & Naive & Sample splitting\\
\midrule
50 & 0.917 & 0.941 & 0.598 & 0.952\\
100 & 0.944 & 0.957 & 0.619 & 0.928\\
400 & 0.939 & 0.943 & 0.564 & 0.944\\
\bottomrule
\end{tabular}
\end{table}


\section{Discussion}

Nothing here is new; the point is that the numbers above were produced by a
pipeline in which every figure and table is tied to the code and the parameter
settings that generated it.
